\chapter{小参数中的计算机证明，实际证明了什么}\label{app:small}
\section{先列出缺口，避免把小例子当作全部证明}
五平均只需额外处理十六、十七、十八元；七平均需要二十二至二十六元。其他位置数已经由一般范围覆盖。

其中五平均的三个点与七平均的二十三至二十六元共用一条路线：假设所有安全候选都发生碰撞，把碰撞写成线性方程，再证明这些方程与零和、同余见证不相容。七平均二十二元另用矩阵群和一个带尺度标签的有限图。

因此这里有两种有限证书：一类穷尽线性方程\emph{模式}，另一类穷尽一个已经证明足够的有限\emph{状态空间}。两者都不对实际输入数值规定高度上限。

\section{一条碰撞等式，怎样变成线性方程}
设两个锚点为 $a,b$，令 $\Delta=a-b$。把每个轻值写成
\[x_i=b+\Delta t_i.
\]
这样 $b$ 对应参数零，$a$ 对应参数一。选择一个轻值子集 $S$，再加 $i$ 份 $a$、若干份 $b$ 补足 $p$ 项，均值为
\[b+\frac{\Delta}{p}\left(i+\sum_{j\in S}t_j\right).
\]
若它恰好碰撞到 $a$，就必须有 $\sum_{j\in S}t_j=p-i$。这是关于实际有理参数的严格线性方程，不只是模 $p$ 同余。

实际 $t_i$ 有无穷多个，但它们的余数 $z_i=t_i\pmod p$ 只有有限种。每个余数模式决定哪些候选和必须写哪些方程。于是可以对有限种模式逐一做精确消元，而消元得到的是每个系统的\emph{全部有理解}。

\section{保护集中的每个位置，为什么有一个私人素数}
先为危险支持各留一个位置，然后删掉所有多余的保护位置，得到一个“再删任何一项都不行”的集合 $T$。对其中某个位置 $v$，总有一个危险素数的例外支持只由它保护；否则 $v$ 还能删。这就是私人素数。

在这个素数处，其他保护位置与锚点 $b$ 同余，只有 $v$ 不同。因此，如果未保护部分相对 $b$ 的总和为零，整个配置的相对总和就只剩一个非零项，违背全局零和。这个矛盾只需要知道\emph{总和}，不必强迫每个轻值都与 $b$ 同余。
\begin{lemma}[五平均和七平均的碰撞分类]\label{lem:collision-small}
五平均的 $n=16,17,18$ 及七平均的 $n=23,24,25,26$ 均满足判据。
\end{lemma}
\begin{proof}
引理\ref{lem:initialize}适用。引理\ref{lem:background}、\ref{lem:heavy}、\ref{lem:two-heavy}只留下不同余的两个重值 $a,b$，$f(a)=p+\alpha,f(b)=p+\beta$、$\alpha\ge\beta\ge0$、$\alpha+\beta\le p-3$。
五平均十六元另有三值 $(a^5,b^6,c^5)$，其中 $a<b<c$，三者模五不同。
选整数调节 $(a^i,c^{4-i},b)$，保留五份 $b$。
若均值碰撞到 $b$，$i=1$ 给 $a+3c=4b$，代入零和得 $13b-5c=0$，
从而 $b\equiv c\pmod2$，再由前式得到 $a\equiv c\pmod2$，违反合法性。
$i=3$ 对称；$i=2$ 给 $a+c=2b$，零和则给 $16b=0$，故 $b=0,a=-c$。
因此这项容量边界也已经处理。

\par\noindent\textbf{保护集与私人素数。}
从需要保护的轻值危险集选一个包含极小横截集 $T$，$h=|T|\le d_p(n)$。
这里“包含极小”是指删去其中任意一个位置就不再同时击中全部危险支持。
因此对每个被保护位置，总有某个危险支持只在这个位置与保护集相交；
否则可以把该位置删掉。这就是所谓私人素数的来源。
每个 $v\in T$ 有私人素数 $\ell$，其危险支持与 $T$ 恰交于 $v$；此处 $a=b\pmod\ell$，其余保护位置与 $b$ 同余。
故未保护部分相对 $b$ 的总和不可能为 $0\pmod\ell$。
若它为零，总和
\[
\sum_i(x_i-b)=\sum_{i\notin T}(x_i-b)+\sum_{i\in T}(x_i-b)
\]
模 $\ell$ 就恰剩私人位置 $v$ 的一个非零差。
但左侧因原配置零和且 $\ell\mid n$ 应为零，矛盾。
这个论证只需要未保护部分的\emph{总和}为零，不要求其中每个值都同余于 $b$。

\par\noindent\textbf{把轻值写成相对锚点的参数。}
令 $\Delta=a-b$、$t_i=(x_i-b)/\Delta$、$z_i=t_i\pmod p$。锚点同余修复后，余数 $0,1$ 各至多一个，实际参数为 $p,1-p$；各余类重数小于 $p$。对轻值子集 $S$，$k=|S|\le p$，若 $i\equiv-\sum_{S}z\pmod p$ 满足 $i\le\min(\alpha,p-k)$，
所选组由轻值子集 $S$、$i$ 份 $a$ 和 $p-k-i$ 份 $b$ 组成，
它的和为 $pb+\Delta(i+\sum_S t)$。
均值碰撞到 $a=b+\Delta$，恰好要求 $i+\sum_S t=p$。
因此碰撞要求
\begin{equation}\label{eq:collision}
\sum_{S}t=p-i;
\end{equation}
若还满足 $p-k-i\le\beta$，两个旧重值均保留，因而该模式不可能失败。从另一锚点，$j\equiv\sum_S z-k\pmod p$、$j\le\min(\beta,p-k)$ 时，碰撞要求 $\sum_S t=k+j-p$。

\par\noindent\textbf{有限线性分类与总和矛盾。}
每个轻值的模 $p$ 余数只有 $p$ 种，各余类的重数也有明确上界。
因此方程模式只有有限个，尽管实际轻值 $t_i$ 可以是任意有理数。
命题\ref{prop:collision-certificate}穷尽这些线性系统。五平均令 $c=n-10-h\in\{5,6,7,8\}$。$c=7,8$ 无解；$c=6$ 的两个解均为整数且 $D=5+\alpha+\sum t_i=15$；$c=5$ 时所有 $D$ 的分母为奇数。对 $n=17,18$，有保护时私人素数整除 $\Delta$，而 $D=15$ 给总和矛盾；无保护时没有相容模式。对 $n=16,h=0$，$16b+15\Delta=0$ 迫使 $\Delta$ 偶，全部整数参数使配置同奇偶；$h=1$ 时 $\Delta$ 偶而 $D$ 分母奇，同样违反私人素数 $2$ 的总和条件。
七平均可在每个余数类内对\eqref{eq:collision}取对称平均，压缩为变量 $y_r$、重数 $f_r$；总和保持。某类在一个方程中部分选取时，交换位置进一步强迫该类实际常值。令 $n-h=21+e$，则 $1\le e\le5$。证书给十一种相容模式，
每个实际参数均为整数，且 $0<D=7+\alpha+\sum f_ry_r<21+e$。$h>0$ 时私人素数给总和矛盾；$h=0$ 时所有值属于 $b+\Delta\Z$，合法性给 $\gcd(\Delta,n)=1$，零和迫使 $n\mid D$，又矛盾。因此总有安全续步，整数能量下降完成。
\end{proof}

\section{如何读一行碰撞表}
五平均表的一行是 $c=6,\alpha=1,\beta=0$，余数列 $(3,4,4,4,4)$，唯一实际解为 $(-7,4,4,4,4)$。这些有理未知数一旦全部满足碰撞方程，就被迫取这五个值。加上锚点贡献，
\[D=5+1+(-7+4+4+4+4)=15.
\]
当位置数为十七或十八、存在保护位置时，私人素数整除 $\Delta$，未保护相对总和 $\Delta D$ 便在该素数处为零，正好触发上面的矛盾。

这不是验证了五个具体输入，而是证明一种可能的失败模式只能有这个解，然后把它排除。
\begin{proposition}[碰撞分类证书]\label{prop:collision-certificate}
在引理\ref{lem:collision-small}的方程下，五平均的903个无序模式恰有表\ref{tab:five-data}所列22个相容解；七平均的57026个压缩模式恰有表\ref{tab:seven-data}所列11个相容解。每个相容系统满秩。
\end{proposition}
\begin{proof}
五平均遍历 $(\alpha,\beta)=(0,0),(1,0),(1,1),(2,0)$、$c=5,6,7,8$，以及长度 $c-\alpha-\beta$ 的非降余数列。删除含五个相同余数、或余数 $0,1$ 各出现超过一次的列。对每个非空至多五项子集，写入\eqref{eq:collision}及另一锚点方程；同时保留两个旧重值时写入矛盾方程 $0=1$。按最左非零元逐列消元，每次整数行消元后约去最大公因子，最后以有理数回代。所得全部相容解如下表。
七平均遍历 $1\le e\le5$、$0\le\beta\le\alpha$、$\alpha+\beta\le4$，长度为 $7+e-\alpha-\beta$。每类重数小于七，余数 $0,1$ 各至多一次。按各类选择数 $0\le k_r\le f_r$、$1\le\sum k_r\le7$ 建立压缩方程，用同一消元。每个重复类都出现在至少一条部分选取方程中，故其实际值等于表列解。程序逐行回代并比较整个解集，得到所述计数及完整表。
\end{proof}
\begingroup\small
\begin{longtable}{rrrlll}
\caption{五平均的全部相容系统}\label{tab:five-data}\\
\toprule $c$&$\alpha$&$\beta$&余数列&唯一解&$D$\\\midrule\endfirsthead
\toprule $c$&$\alpha$&$\beta$&余数列&唯一解&$D$\\\midrule\endhead
5&0&0&$(0,2,2,2,2)$&$(5,-8,-8,-8,-8)$&$-22$\\
5&0&0&$(0,3,3,3,3)$&$(5,-2,-2,-2,-2)$&$2$\\
5&0&0&$(0,4,4,4,4)$&$(5,\frac{-7}{2},\frac{-7}{2},\frac{-7}{2},\frac{-7}{2})$&$-4$\\
5&0&0&$(1,2,2,2,2)$&$(-4,\frac{9}{2},\frac{9}{2},\frac{9}{2},\frac{9}{2})$&$19$\\
5&0&0&$(1,3,3,3,3)$&$(-4,3,3,3,3)$&$13$\\
5&0&0&$(1,4,4,4,4)$&$(-4,9,9,9,9)$&$37$\\
5&0&0&$(2,2,2,2,3)$&$(-3,-3,-3,-3,8)$&$1$\\
5&0&0&$(2,2,2,2,4)$&$(7,7,7,7,-16)$&$17$\\
5&0&0&$(2,3,3,3,3)$&$(12,-7,-7,-7,-7)$&$-11$\\
5&0&0&$(2,4,4,4,4)$&$(17,-6,-6,-6,-6)$&$-2$\\
5&0&0&$(3,3,3,3,4)$&$(8,8,8,8,-11)$&$26$\\
5&0&0&$(3,4,4,4,4)$&$(-7,4,4,4,4)$&$14$\\
5&1&0&$(2,2,2,2)$&$(2,2,2,2)$&$14$\\
5&1&0&$(3,3,3,3)$&$(\frac{4}{3},\frac{4}{3},\frac{4}{3},\frac{4}{3})$&$\frac{34}{3}$\\
5&1&0&$(3,4,4,4)$&$(-7,4,4,4)$&$11$\\
5&1&0&$(4,4,4,4)$&$(4,4,4,4)$&$22$\\
6&1&0&$(3,4,4,4,4)$&$(-7,4,4,4,4)$&$15$\\
5&1&1&$(0,3,3)$&$(5,-7,-7)$&$-3$\\
5&1&1&$(1,3,3)$&$(-4,8,8)$&$18$\\
5&2&0&$(2,2,2)$&$(2,2,2)$&$13$\\
5&2&0&$(2,2,3)$&$(2,2,3)$&$14$\\
6&2&0&$(2,2,2,2)$&$(2,2,2,2)$&$15$\\
\bottomrule\end{longtable}\endgroup

\begingroup\small
\begin{longtable}{rrrllll}
\caption{七平均的全部相容系统}\label{tab:seven-data}\\
\toprule $e$&$\alpha$&$\beta$&余数&重数&唯一解&$D$\\\midrule\endfirsthead
\toprule $e$&$\alpha$&$\beta$&余数&重数&唯一解&$D$\\\midrule\endhead
1&2&0&$(2)$&$(6)$&$(2)$&$21$\\
1&3&0&$(2)$&$(5)$&$(2)$&$20$\\
1&3&0&$(2,3)$&$(4,1)$&$(2,3)$&$21$\\
1&4&0&$(2)$&$(4)$&$(2)$&$19$\\
1&4&0&$(2,3)$&$(3,1)$&$(2,3)$&$20$\\
1&4&0&$(2,4)$&$(3,1)$&$(2,4)$&$21$\\
1&4&0&$(2,3)$&$(2,2)$&$(2,3)$&$21$\\
2&3&0&$(2)$&$(6)$&$(2)$&$22$\\
2&4&0&$(2)$&$(5)$&$(2)$&$21$\\
2&4&0&$(2,3)$&$(4,1)$&$(2,3)$&$22$\\
3&4&0&$(2)$&$(6)$&$(2)$&$23$\\
\bottomrule\end{longtable}\endgroup

\section{二十二元的3872个状态是什么}
控制模数为 $M=1936$。列记为 $(z,y)$，其中 $y$ 是模 $M$ 单位。先记录比值 $x=z/y\pmod M$，有1936种。

但实际可用的环 $R=\Z[1/35]$ 的单位模 $M$ 只覆盖全部单位的一半，记这半为 $U$，另一半为 $3U$。所以再加一位 $\epsilon$：$y\in U$ 时为零，$y\in3U$ 时为一。总共 $1936\times2=3872$ 个状态。

这是为何单纯遍历1936个方向不够。两个方向相同的列，可能一个处于允许单位倍数类，另一个不在；必须看真实宏能否把它们之间的尺度差也调整掉。
\begin{lemma}[七平均二十二元的方向与尺度]\label{lem:seven22}
七平均二十二元的合法配置均可归零。
\end{lemma}
\begin{proof}

\par\noindent\textbf{第一步：三个真实返回与深层。}
入口使用命题\ref{prop:entry}证明中转换前的 $(a^{14},(-2a+z)^7,-7z)$。三条三步返回为
\[
A=\begin{pmatrix}4&1\\-1&0\end{pmatrix},\quad
F=\begin{pmatrix}1&2\\0&7\end{pmatrix},\quad
J=\begin{pmatrix}1&2\\2&-1\end{pmatrix}.
\]
$A$ 保留一份旧 $A$、平均两组 $(A^6,B)$ 再平均 $(A,B^5,C)$；$F$ 保留 $C$、平均两组 $(A^5,B^2)$ 再平均 $(A^4,B^3)$；$J$ 保留一份旧 $B$、平均两组 $(A^5,B^2)$ 再平均 $(A^4,B^2,C)$。实际尺度均为 $1/7$。有 $(FA^2)^2=-7I,J^2=5I$，故各逆可实现。
令 $H=AFJA^{-1}=\left(\begin{smallmatrix}-7&-62\\0&5\end{smallmatrix}\right)$，则 $[F,H]=U(396/35)$。倍率生成 $R=\Z[1/35]$；用 $C=\left(\begin{smallmatrix}1&4\\0&-1\end{smallmatrix}\right)$ 换基，$CU(t)C=U(-t)$、$CAU(t)A^{-1}C=L(t)$，得到 $\Gamma(396^2,R)$。

\par\noindent\textbf{第二步：去掉多余的三进控制层。}
记 $M=1936$、$X=F^2/7$。命题\ref{prop:seven-certificate}给 $A^{440}=X^{440}=I\pmod M$，且两者模 $9$ 生成全部 $\SL_2(\Z/9\Z)$。其中三个方向 $I+3E_{12},I+3E_{21},I+3\diag(1,-1)$ 的立方逐层生成模 $81$ 的高层核。因此对任意 $g\in\Gamma(M,R)$，可选可实现 $w\equiv I\pmod M$、$w\equiv g\pmod{81}$；由 $396^2=81M$ 得 $w^{-1}g$ 已在深核中，故 $\Gamma(M,R)$ 可实现。

\par\noindent\textbf{第三步：有限方向与尺度同时命中终端。}
同一证书给 $R^\times$ 模 $M$ 的像 $U$ 有 $440$ 个元素，全部单位为 $U\sqcup3U$。合法列 $(z,y)=(z,a+7z)$ 的有限标记是 $(z/y,\epsilon)$，$\epsilon$ 记录 $y$ 所属单位陪集。证书证明全部 $3872$ 个标记都能到 $(0,0)$。引理\ref{lem:fiber}再精确送到 $(0,s)$，其中 $s\in R^\times$。配置 $(s^{14},(-2s)^7,0)$ 中，四份 $s$、两份 $-2s$ 和一个零组成零和七元组，接引理\ref{lem:zero-tail}。
\end{proof}

\section{有限图的边，怎样提升成无限输入的操作}
矩阵 $(a,b;c,d)$ 作用于列后，第二坐标乘上 $cx+d$。于是 $x$ 按分式线性变换更新，而 $\epsilon$ 要加上 $cx+d$ 属于哪一类的标记。

证书从终端 $(0,0)$ 出发，把每条允许边都展开，直到没有新状态；结果恰好覆盖3872个。生成元已经有真实逆，所以这条父边树可以反向走回终端。至于实际整数有多大，最后由 $\Gamma(M,R)$ 的精确纤维引理处理。
\begin{proposition}[二十二元有限轨道证书]\label{prop:seven-certificate}
引理\ref{lem:seven22}中的 $A^{440},(F^2/7)^{440}$ 模 $1936$ 均为恒等，模 $9$ 生成 $648$ 个行列式一矩阵。环 $\Z[1/35]$ 的单位像模 $1936$ 为 $440$ 元子群 $U$，单位群为 $U\sqcup3U$。三个返回及其逆在 $1936\times2$ 个方向--尺度标记上作用传递，最大父边深度为八。
\end{proposition}
\begin{proof}
快速整数模幂检查前一等式。从恒等矩阵出发，逐个乘两个模九生成元，保留尚未出现的矩阵，闭包大小为648；逐项检查行列式为一。因 $|\SL_2(\Z/9\Z)|=3^3(3^2-1)3=648$，这就是全群。
同理从 $1$ 乘 $-1,5,7$ 得 $U$，并检查 $U,3U$ 不交且并为全部880个单位。令 $\chi$ 为其二元商标记。在 $(z,y)$ 坐标中使用
\[
A=\begin{pmatrix}7&-1\\22&-3\end{pmatrix},\quad
F=\begin{pmatrix}7&0\\44&1\end{pmatrix},\quad
J=\begin{pmatrix}-15&2\\-110&15\end{pmatrix}.
\]
矩阵 $(a,b;c,d)$ 对标记的作用是
\[
(x,\epsilon)\longmapsto((ax+b)(cx+d)^{-1},\epsilon+\chi(cx+d)).
\]
从 $(0,0)$ 对 $A,A^{-1},F,F^{-1},J$ 逐层闭包，得到全部3872个状态，最深层为八。对每个状态的全部边重新检查分母可逆及目标在集合内，即给完整有限证书。
\end{proof}

“最大父边深度为八”只说有限标记图最多走八条这样的边，不能直接说原题只需八次平均。每条宏本身有若干原子操作，精确纤维输送也还需要平均词。有限图深度与原子词长度是两个不同计数。
